\documentclass[xcolor=pdftex,dvipsnames,table]{beamer}
\mode<presentation>
\usetheme{boxes}
\setbeamertemplate{navigation symbols}{}
% http://www.latex-community.org/forum/viewtopic.php?f=4&t=6694
\setbeamertemplate{navigation symbols}{\raisebox{5pt}{\makebox[\paperwidth]{\hfill\makebox[10pt]{\scriptsize\insertframenumber\vspace{1ex}}}}}
\setbeamertemplate{footline}[frame number]
\setbeamertemplate{blocks}[shadow=false]
\setbeamercolor*{block title}{fg=structure,bg=RoyalBlue!10}
\setbeamercolor*{block title example}{fg=BrickRed,bg=Goldenrod!10}
\setbeamercolor*{block title alerted}{fg=white,bg=black}
\addtobeamertemplate{block begin}{\pgfsetfillopacity{0.8}}{\pgfsetfillopacity{1}}
\rowcolors{1}{RoyalBlue!20}{RoyalBlue!5}

%\DeclareGraphicsRule{*}{mps}{*}{}

\usepackage{latexsym}
\usepackage{hyperref}

\raggedright

\newcount\lecturecount
\lecturecount=0
\AtBeginLecture{%
    \advance\lecturecount by 1
    \date{}
    \begin{frame}
    \begin{center}
    \titlepage
    \ifnum\lecturecount=1
    Part \the\lecturecount: \insertlecture
    \else
    Part \the\lecturecount: \insertlecture
    \fi
    \end{center}
    \end{frame}
}

\begin{document}


\title{\color{blue}Natural Language Processing}

\author{Anoop Sarkar \\ {\tt http://anoopsarkar.github.io/nlp-class}}
\institute{}
%\date{}
     
{
\setbeamertemplate{navigation symbols}{}
\addtocounter{framenumber}{-1}
\begin{frame}
\begin{center}
\vspace{1cm}
\includegraphics[scale=0.4]{figures/natlang-cky-logo.eps}
\end{center}
\titlepage
\end{frame}
}



\lecture{Generative Models for Word Alignment}{}

\section{Statistical Machine Translation}
\frame{\tableofcontents[currentsection]}

\begin{frame}
\frametitle{Statistical Machine Translation}
\begin{block}{Noisy Channel Model}
\begin{eqnarray*}
\textbf{e}^\ast & = & \arg\max_{\textbf{e}} \underbrace{\Pr(\textbf{e})}_{\color{red}\textbf{Language Model}} \pause \cdot \underbrace{\Pr(\textbf{f} \mid \textbf{e})}_{\color{blue}\textbf{Alignment Model}}
\end{eqnarray*}
\end{block}
\end{frame}

\begin{frame}
\def\blockdist{2.0}
\begin{alertblock}{Alignment Task}
\begin{tikzpicture}[node distance=2.5cm,auto,>=latex']
    \node (a) {Program};
    \path (a.140)+(-\blockdist,0) node (b) {\textbf{e}};
    \path (a.-140)+(-\blockdist,0) node (c) {\textbf{f}};
    \node (d) [right of=a] {$\Pr(\textbf{e} \mid \textbf{f})$};
    \path[->] (b) edge (a);
    \path[->] (c) edge (a);
    \path[->] (a) edge (d);
\end{tikzpicture}
\end{alertblock}\pause
\begin{alertblock}{Training Data}
\begin{itemize}[<+->]
\item \color{blue}\textbf{Alignment Model}: learn a mapping between \textbf{f} and \textbf{e}. \\
Training data: lots of translation pairs between \textbf{f} and \textbf{e}.
\end{itemize}
\end{alertblock}
\end{frame}

\begin{frame}
\frametitle{Statistical Machine Translation}
\begin{block}{The IBM Models}
\begin{itemize}[<+->]
\item The first statistical machine translation models were developed at IBM Research (Yorktown Heights, NY) in the 1980s
\item The models were published in 1993: \\
{\small Brown et.\ al.\ The Mathematics of Statistical Machine Translation. \textit{Computational Linguistics}. 1993.} \\
{\small \url{http://aclweb.org/anthology/J/J93/J93-2003.pdf}}
\item These models are the basic SMT models, called: \\
IBM Model 1, IBM Model 2, IBM Model 3, IBM Model 4, IBM Model 5 \\
as they were called in the 1993 paper.
\item We use \textbf{e} and \textbf{f} in the equations in honor of their system which translated from French to English.\\
Trained on the Canadian Hansards (Parliament Proceedings)
%\item Twenty years later, they came together at a workshop called {\color{blue}\hyperlink{https://sites.google.com/site/20yearsofbitext/}{Twenty Years of Bitext}} \\
%Transcript: \url{http://cs.jhu.edu/~post/bitext/}
\end{itemize}
\end{block}
\end{frame}


\section{Generative Model of Word Alignment}
\frame{\tableofcontents[currentsection]}

\begin{frame}
\frametitle{Generative Model of Word Alignment}
\begin{block}{}
\begin{itemize}
\item English \textbf{e}: Mary did not slap the green witch
\item ``French'' \textbf{f}: Maria no daba una botefada a la bruja verde
\item Alignment \textbf{a}: $\{ 1, 3, 4, 4, 4, 5, 5, 7, 6 \}$ \\
e.g.\ $(f_8, e_{a_8})$ = $(f_8, e_7)$ = (bruja, witch)
\end{itemize}
\end{block}\pause
\begin{block}{Visualizing alignment \textbf{a}}
\begin{tikzpicture}
\begin{scope}[start chain=going right,node distance=2mm]
\def\ephr{0/{Mary},1/{did},2/{not},3/{slap},4/{the},5/{green},6/{witch}}
\foreach \i / \element in \ephr
  \node [on chain] (e\i) {\strut \element};
\end{scope}

\begin{scope}[yshift=-16mm,start chain=going right,node distance=2mm]
\def\sphr{0/{Maria},1/{no},2/{daba},3/{una},4/{botefada},5/{a},6/{la},7/{bruja},8/{verde}}
\foreach \j / \element in \sphr
  \node [on chain] (s\j) {\strut \element};
\end{scope}

\draw (e0) -- (s0);
\draw (e2) -- (s1);
\draw (e3) -- (s2);
\draw (e3) -- (s3);
\draw (e3) -- (s4);
\draw (e4) -- (s5);
\draw (e4) -- (s6);
\draw (e5) -- (s8);
\draw (e6) -- (s7);
\end{tikzpicture}

\end{block}
\end{frame}

\begin{frame}
\frametitle{Generative Model of Word Alignment}
\begin{block}{Data Set}
\begin{itemize}[<+->]
\item Data set ${\cal D}$ of $N$ sentences:
\[ {\cal D} = \{ (\textbf{f}^{(1)}, \textbf{e}^{(1)}), \ldots, (\textbf{f}^{(N)}, \textbf{e}^{(N)}) \} \]
\item French $\textbf{f}$: $( f_1, f_2, \ldots, f_I )$
\item English $\textbf{e}$: $( e_1, e_2, \ldots, e_J )$
\item Alignment $\textbf{a}$: $( a_1, a_2, \ldots, a_I )$
\end{itemize}
\end{block}
\end{frame}

\begin{frame}
\frametitle{Generative Model of Word Alignment}
\begin{block}{Find the best alignment for each translation pair}
\[ \textbf{a}^\ast = \arg\max_{\textbf{a}} \Pr(\textbf{a} \mid \textbf{f}, \textbf{e}) \]
\end{block}\pause
\begin{block}{Chain rule revisited}
\begin{eqnarray*}
\Pr( \textbf{f}, \textbf{a} ) & = & \Pr(\textbf{f}) \Pr(\textbf{a} \mid \textbf{f}) \pause \\
\Pr( \textbf{f}, \textbf{a} \mid \textbf{e} ) & = & \Pr(\textbf{f} \mid \textbf{e}) \Pr(\textbf{a} \mid \textbf{f}, \textbf{e}) 
\end{eqnarray*}
\end{block}
\begin{block}{Alignment probability}
\[ \Pr(\textbf{a} \mid \textbf{f}, \textbf{e}) = \frac{ \Pr( \textbf{f}, \textbf{a} \mid \textbf{e} ) }{ \Pr(\textbf{f} \mid \textbf{e}) } = \frac{ \Pr( \textbf{f}, \textbf{a} \mid \textbf{e} ) }{ \sum_{\textbf{a}} \Pr(\textbf{f}, \textbf{a} \mid \textbf{e}) } 
\]
\end{block}
\end{frame}

\subsection{Word Alignments: IBM Model 3}
\frame{\tableofcontents[currentsection]}

\begin{frame}
\frametitle{Word Alignments: IBM Model 3}
\begin{block}{Generative ``story'' for $P(\textbf{a} \mid \textbf{f}, \textbf{e})$}
\begin{tikzpicture}
\begin{scope}[start chain=going right,node distance=2mm]
\def\fphr{0/{Mary},1/{did},2/{not},3/{slap},4/{the},5/{green},6/{witch}}
\foreach \i / \element in \fphr
  \node [on chain] (k\i) {\strut \element};
\end{scope}

\begin{scope}[yshift=-16mm,start chain=going right,node distance=2mm]
\def\ephr{0/{Mary},1/{not},2/{slap},3/{slap},4/{slap},5/{the},6/{the},7/{green},8/{witch}}
\foreach \j / \element in \ephr
  \node [on chain] (e\j) {\strut \element};
\end{scope}

\draw[->] (k0) -- (e0);
\draw[->] (k2) -- (e1);
\draw[->] (k3) -- (e2);
\draw[->] (k3) -- (e3);
\draw[->] (k3) -- (e4);
\draw[->] (k4) -- (e5);
\draw[->] (k4) -- (e6);
\draw[->] (k5) -- (e7);
\draw[->] (k6) -- (e8);

\pause 

\begin{scope}[yshift=-32mm,start chain=going right,node distance=2mm]
\def\sphr{0/{Maria},1/{no},2/{daba},3/{una},4/{botefada},5/{a},6/{la},7/{verde},8/{bruja}}
\foreach \j / \element in \sphr
  \node [on chain] (s\j) {\strut \element};
\end{scope}

\draw[->] (e0) -- (s0);
\draw[->] (e1) -- (s1);
\draw[->] (e2) -- (s2);
\draw[->] (e3) -- (s3);
\draw[->] (e4) -- (s4);
\draw[->] (e5) -- (s5);
\draw[->] (e6) -- (s6);
\draw[->] (e7) -- (s7);
\draw[->] (e8) -- (s8);

\pause

\begin{scope}[yshift=-48mm,start chain=going right,node distance=2mm]
\def\rphr{0/{Maria},1/{no},2/{daba},3/{una},4/{botefada},5/{a},6/{la},7/{bruja},8/{verde}}
\foreach \j / \element in \rphr
  \node [on chain] (r\j) {\strut \element};
\end{scope}

\draw[->] (s0) -- (r0);
\draw[->] (s1) -- (r1);
\draw[->] (s2) -- (r2);
\draw[->] (s3) -- (r3);
\draw[->] (s4) -- (r4);
\draw[->] (s5) -- (r5);
\draw[->] (s6) -- (r6);
\draw[->] (s7) -- (r8);
\draw[->] (s8) -- (r7);

\end{tikzpicture}
\end{block}
\end{frame}

\begin{frame}
\frametitle{Word Alignments: IBM Model 3}
\begin{block}{Fertility parameter}
\[ n(\phi_j \mid e_j) : n(3 \mid \textit{slap}) \]
\end{block}\pause
\begin{block}{Translation parameter}
\[ t(f_i \mid e_j) : t(\textit{bruja} \mid \textit{witch}) \]
\end{block}\pause
\begin{block}{Distortion parameter}
\[ d(f_{\textit{pos}} \mid e_{\textit{pos}}, I, J) : d(8 \mid 7, 8, 6) \]
\end{block}
\end{frame}

\subsection{Word Alignments: IBM Model 1}
\frame{\tableofcontents[currentsection]}

\begin{frame}
\frametitle{Word Alignments: IBM Model 1}
\begin{block}{Alignment probability}
\[ \Pr(\textbf{a} \mid \textbf{f}, \textbf{e}) = \frac{ \Pr( \textbf{f}, \textbf{a} \mid \textbf{e} ) }{ \sum_{\textbf{a}} \Pr(\textbf{f}, \textbf{a} \mid \textbf{e}) } 
\]
\end{block}\pause
\begin{columns}[t]

\begin{column}{\textwidth/2}
\begin{block}{Example alignment}
\begin{tikzpicture}
\begin{scope}[start chain=going right,node distance=2mm]
\def\dphr{1/{das},2/{Haus},3/{ist},4/{klein}}
\foreach \i / \element in \dphr
  \node [on chain] (d\i) {\strut $\stackrel{\i}{\mbox{\element}}$};
\end{scope}

\begin{scope}[yshift=-14mm,start chain=going right,node distance=2mm]
\def\ephr{1/{the},2/{house},3/{is},4/{small}}
\foreach \i / \element in \ephr
  \node [on chain] (e\i) {\strut $\stackrel{\i}{\mbox{\element}}$};
\end{scope}

\draw (d1) -- (e1);
\draw (d2) -- (e2);
\draw (d3) -- (e3);
\draw (d4) -- (e4);

\end{tikzpicture}
\end{block}
\end{column}\pause

\begin{column}{\textwidth/2}
\begin{block}{$\Pr(\textbf{f}, \textbf{a} \mid \textbf{e}) = \prod_{i=1}^I t(f_i \mid e_{a_i})$}
\begin{eqnarray*}
\lefteqn{\Pr( \textbf{f}, \textbf{a} \mid \textbf{e} ) =}\\
&& t(\textrm{das} \mid \textrm{the}) \times \\
&& t(\textrm{Haus} \mid \textrm{house}) \times \\
&& t(\textrm{ist} \mid \textrm{is}) \times \\
&& t(\textrm{klein} \mid \textrm{small})
\end{eqnarray*}
\end{block}
\end{column}

\end{columns}
\end{frame}


\begin{frame}
\frametitle{Word Alignments: IBM Model 1}
\begin{block}{Sum over all alignments}
\begin{eqnarray*}
\sum_{\textbf{a}} \Pr(\textbf{f}, \textbf{a} \mid \textbf{e}) & = & 
\sum_{a_1 = 1}^J \sum_{a_2 = 1}^J \ldots \sum_{a_I = 1}^J \prod_{i=1}^I t(f_i \mid e_{a_i}) 
\end{eqnarray*}
\end{block}\pause
\begin{block}{Assume $(f_1, f_2, f_3)$ and $(e_1, e_2)$}
\[ \sum_{a_1=1}^2 \sum_{a_2=1}^2 \sum_{a_3=1}^2 t(f_1 \mid e_{a_1}) \times t(f_2 \mid e_{a_2}) \times t(f_3 \mid e_{a_3}) \]
\end{block}
\end{frame}

\begin{frame}
\frametitle{Word Alignments: IBM Model 1}
\begin{block}{Assume $(f_1, f_2, f_3)$ and $(e_1, e_2)$: $I=3$ and $J=2$}
\[ \sum_{a_1=1}^2 \sum_{a_2=1}^2 \sum_{a_3=1}^2 t(f_1 \mid e_{a_1}) \times t(f_2 \mid e_{a_2}) \times t(f_3 \mid e_{a_3}) \]
\end{block}
\begin{block}{$J^I = 2^3$ terms to be added:}
\[ 
\begin{array}{cccccc}
\color{red} t(f_1 \mid e_1) & \times & \color{red} t(f_2 \mid e_1) & \times & \color{red} t(f_3 \mid e_1) & + \\
\color{red} t(f_1 \mid e_1) & \times & \color{red} t(f_2 \mid e_1) & \times & \color{red} t(f_3 \mid e_2) & + \\
\color{blue} t(f_1 \mid e_1) & \times & \color{blue} t(f_2 \mid e_2) & \times & \color{blue} t(f_3 \mid e_1) & + \\
\color{blue} t(f_1 \mid e_1) & \times & \color{blue} t(f_2 \mid e_2) & \times & \color{blue} t(f_3 \mid e_2) & + \\
\color{red} t(f_1 \mid e_2) & \times & \color{red} t(f_2 \mid e_1) & \times & \color{red} t(f_3 \mid e_1) & + \\
\color{red} t(f_1 \mid e_2) & \times & \color{red} t(f_2 \mid e_1) & \times & \color{red} t(f_3 \mid e_2) & + \\
\color{blue} t(f_1 \mid e_2) & \times & \color{blue} t(f_2 \mid e_2) & \times & \color{blue} t(f_3 \mid e_1) & + \\
\color{blue} t(f_1 \mid e_2) & \times & \color{blue} t(f_2 \mid e_2) & \times & \color{blue} t(f_3 \mid e_2) & + \\
\end{array}
\]
\end{block}
\end{frame}

\begin{frame}
\frametitle{Word Alignments: IBM Model 1}
\begin{block}{Factor the terms:}
\[ 
\begin{array}{cccc}
( t(f_1 \mid e_1) \times t(f_2 \mid e_1) ) & \times & ( t(f_3 \mid e_1) + t(f_3 \mid e_2) ) & + \\
( t(f_1 \mid e_1) \times t(f_2 \mid e_2) ) & \times & ( t(f_3 \mid e_1) + t(f_3 \mid e_2) ) & + \\
( t(f_1 \mid e_2) \times t(f_2 \mid e_1) ) & \times & ( t(f_3 \mid e_1) + t(f_3 \mid e_2) ) & + \\
( t(f_1 \mid e_2) \times t(f_2 \mid e_2) ) & \times & ( t(f_3 \mid e_1) + t(f_3 \mid e_2) ) & 
\end{array}
\]
\pause
\[ \left( t(f_3 \mid e_1) + t(f_3 \mid e_2) \right) \left(
\begin{array}{cccc}
 t(f_1 \mid e_1) & \times & t(f_2 \mid e_1)  & + \\
 t(f_1 \mid e_1) & \times & t(f_2 \mid e_2)  & + \\
 t(f_1 \mid e_2) & \times & t(f_2 \mid e_1)  & + \\
 t(f_1 \mid e_2) & \times & t(f_2 \mid e_2)  & 
\end{array}
\right) \]
\pause
\[ {\color{RoyalBlue} \left( t(f_3 \mid e_1) + t(f_3 \mid e_2) \right) } \left(
\begin{array}{cccc}
 \color{red} t(f_1 \mid e_1) & \times & \color{blue} ( t(f_2 \mid e_1) + t(f_2 \mid e_2) ) & + \\
 \color{red} t(f_1 \mid e_2) & \times & \color{blue} ( t(f_2 \mid e_1) + t(f_2 \mid e_2) ) &
\end{array}
\right) \]
\end{block}
\end{frame}

\begin{frame}
\frametitle{Word Alignments: IBM Model 1}
\begin{block}{Assume $(f_1, f_2, f_3)$ and $(e_1, e_2)$: $I=3$ and $J=2$}
\[ \prod_{i=1}^3 \sum_{a_i=1}^2 t(f_i \mid e_{a_i}) \]
\end{block}
\begin{block}{$I \times J = 2 \times 3$ terms to be added:}
\[ 
\begin{array}{cccc}
\color{red} ( t(f_1 \mid e_1) & + & \color{red} t(f_1 \mid e_2) ) & \times \\
\color{blue} ( t(f_2 \mid e_1) & + & \color{blue} t(f_2 \mid e_2) ) & \times \\
\color{RoyalBlue} ( t(f_3 \mid e_1) & + & \color{RoyalBlue} t(f_3 \mid e_2) ) & 
\end{array}
\]
\end{block}

\end{frame}


\section*{Acknowledgements}

\begin{frame}
\centering
\begin{alertblock}{Acknowledgements}
Many slides borrowed or inspired from lecture notes by Michael Collins, Chris Dyer, Kevin Knight, Philipp Koehn, Adam Lopez, and Luke Zettlemoyer from their NLP course materials. 

\bigskip

All mistakes are my own.
\end{alertblock}
\end{frame}



\end{document}
